% GENERATED FILE. Edit canonical lessons, then run npm run generate:books.
% canonical-source-hash: fnv1a64:1c7a5fdd07e41dfa
% canonical-lessons: KA-C370-patrikodyama, KA-C370-saksyacitra, KA-C370-oduga, KA-C370-keluga, KA-C370-kukhyata

\chapter{Journalism and Its Public}
\label{ch:ka-journalism-and-its-public}
\hlchaptermodality{370}

\begin{chapteropening}
\textbf{By the end of this chapter:} \emph{I can say journalism, a documentary, a reader, a listener and notorious.}

Complete the last lesson of chapter 370: I can say journalism, a documentary, a reader, a listener and notorious.
\end{chapteropening}

\section[\texorpdfstring{\kn{ಪತ್ರಿಕೋದ್ಯಮ} (patrikōdyama)}{ಪತ್ರಿಕೋದ್ಯಮ (patrikōdyama)}]{\kn{ಪತ್ರಿಕೋದ್ಯಮ} (patrikōdyama) — journalism}

\label{lesson:KA-C370-patrikodyama}



\begin{quote}
Before the new one: say the Kannada for a jail, then the Kannada for to appear.
\end{quote}

\subsection*{You'll want to know: \kn{ಪತ್ರಿಕೋದ್ಯಮ}}
\textbf{\kn{ಪತ್ರಿಕೋದ್ಯಮ}} — \emph{patrikōdyama} — \textquotedblleft{}journalism\textquotedblright{}.

\begin{quote}
\kn{ಅವಳು} \kn{ಕಾಲೇಜಿನಲ್ಲಿ} \kn{ಪತ್ರಿಕೋದ್ಯಮ} \kn{ಕಲಿತಳು}.
\end{quote}

\begin{quote}
\emph{avaḷu kālējinalli patrikōdyama kalitaḷu.}
\end{quote}

\begin{quote}
She studied journalism at college.
\end{quote}

\kn{ಪತ್ರಿಕೆ}, a newspaper, and \kn{ಉದ್ಯಮ}, an enterprise: the newspaper trade.

\begin{cousinweb}
\textbf{\kn{ಪತ್ರಿಕೋದ್ಯಮ}} (\emph{patrikōdyama}) is borrowed from Sanskrit, a modern compound of \emph{patrikā}, \textquotedblleft{}a leaf, a paper\textquotedblright{}, from \emph{pattra} (CDIAL 7733), and \emph{udyama}, \textquotedblleft{}undertaking, striving\textquotedblright{} (CDIAL 2048).
\end{cousinweb}

\begin{grammarlens}[title={What you've built}]
One more word for telling what happened.
\end{grammarlens}

\subsection*{Your turn}
\begin{itemize}
\raggedright
  \item \emph{Say it:} \emph{patrikōdyama}
  \item \emph{Say it:} \emph{patrikōdyama}, once more
  \item \emph{Say it:} \emph{avaḷu kālējinalli patrikōdyama kalitaḷu.} — She studied journalism at college.
  \item \emph{Recall:} say the Kannada for a jail, then the Kannada for to appear, then say \emph{patrikōdyama} again
\end{itemize}

\begin{tcolorbox}[breakable,colback=teal!4,colframe=teal!35!black,title={Before you move on}]
What does \kn{ಪತ್ರಿಕೋದ್ಯಮ} mean? (\textquotedblleft{}Journalism\textquotedblright{}.) Say it once more.
\end{tcolorbox}
\section[\texorpdfstring{\kn{ಸಾಕ್ಷ್ಯಚಿತ್ರ} (sākṣyacitra)}{ಸಾಕ್ಷ್ಯಚಿತ್ರ (sākṣyacitra)}]{\kn{ಸಾಕ್ಷ್ಯಚಿತ್ರ} (sākṣyacitra) — a documentary}

\label{lesson:KA-C370-saksyacitra}



\begin{quote}
Before the new one: say the Kannada for to appear, then the Kannada for journalism.
\end{quote}

\subsection*{You'll want to know: \kn{ಸಾಕ್ಷ್ಯಚಿತ್ರ}}
\textbf{\kn{ಸಾಕ್ಷ್ಯಚಿತ್ರ}} — \emph{sākṣyacitra} — \textquotedblleft{}a documentary\textquotedblright{}.

\begin{quote}
\kn{ಆ} \kn{ಸಾಕ್ಷ್ಯಚಿತ್ರ} \kn{ಕಾಡಿನ} \kn{ಹುಲಿಗಳ} \kn{ಬಗ್ಗೆ} \kn{ಇದೆ}.
\end{quote}

\begin{quote}
\emph{ā sākṣyacitra kāḍina huligaḷa bagge ide.}
\end{quote}

\begin{quote}
That documentary is about the tigers of the forest.
\end{quote}

Plural: \kn{ಸಾಕ್ಷ್ಯಚಿತ್ರಗಳು}, \emph{sākṣyacitragaḷu}.

\kn{ಸಾಕ್ಷ್ಯ}, evidence, and \kn{ಚಿತ್ರ}, a picture: a film of what is real.

\begin{cousinweb}
\textbf{\kn{ಸಾಕ್ಷ್ಯಚಿತ್ರ}} (\emph{sākṣyacitra}) is borrowed from Sanskrit, a modern compound of \emph{sākṣya}, \textquotedblleft{}evidence\textquotedblright{} (CDIAL 13322), and \emph{citra}, \textquotedblleft{}bright, a picture\textquotedblright{} (CDIAL 4803).
\end{cousinweb}

\begin{grammarlens}[title={What you've built}]
One more word for telling what happened.
\end{grammarlens}

\subsection*{Your turn}
\begin{itemize}
\raggedright
  \item \emph{Say it:} \emph{sākṣyacitra}
  \item \emph{Say it:} \emph{sākṣyacitra}, once more
  \item \emph{Say it:} \emph{ā sākṣyacitra kāḍina huligaḷa bagge ide.} — That documentary is about the tigers of the forest.
  \item \emph{Recall:} say the Kannada for to appear, then the Kannada for journalism, then say \emph{sākṣyacitra} again
\end{itemize}

\begin{tcolorbox}[breakable,colback=teal!4,colframe=teal!35!black,title={Before you move on}]
What does \kn{ಸಾಕ್ಷ್ಯಚಿತ್ರ} mean? (\textquotedblleft{}A documentary\textquotedblright{}.) Say it once more.
\end{tcolorbox}
\section[\texorpdfstring{\kn{ಓದುಗ} (ōduga)}{ಓದುಗ (ōduga)}]{\kn{ಓದುಗ} (ōduga) — a reader}

\label{lesson:KA-C370-oduga}



\begin{quote}
Before the new one: say the Kannada for journalism, then the Kannada for a documentary.
\end{quote}

\subsection*{You'll want to know: \kn{ಓದುಗ}}
\textbf{\kn{ಓದುಗ}} — \emph{ōduga} — \textquotedblleft{}a reader\textquotedblright{}.

\begin{quote}
\kn{ಈ} \kn{ಪತ್ರಿಕೆಗೆ} \kn{ಲಕ್ಷ} \kn{ಓದುಗರು} \kn{ಇದ್ದಾರೆ}.
\end{quote}

\begin{quote}
\emph{ī patrikege lakṣa ōdugaru iddāre.}
\end{quote}

\begin{quote}
This newspaper has a hundred thousand readers.
\end{quote}

Plural: \kn{ಓದುಗರು}, \emph{ōdugaru}.

-\kn{ಗ} makes \textquotedblleft{}one who does\textquotedblright{}: \kn{ಓದು}, to read, gives \kn{ಓದುಗ}, a reader.

\begin{cousinweb}
\textbf{\kn{ಓದುಗ}} (\emph{ōduga}) is a native Kannada noun: \emph{ōdu}, \textquotedblleft{}to read, to recite\textquotedblright{} (DEDR 1052), with the ending -\emph{ga}, \textquotedblleft{}one who does\textquotedblright{}.
\end{cousinweb}

\begin{grammarlens}[title={What you've built}]
One more word for telling what happened.
\end{grammarlens}

\subsection*{Your turn}
\begin{itemize}
\raggedright
  \item \emph{Say it:} \emph{ōduga}
  \item \emph{Say it:} \emph{ōduga}, once more
  \item \emph{Say it:} \emph{ī patrikege lakṣa ōdugaru iddāre.} — This newspaper has a hundred thousand readers.
  \item \emph{Recall:} say the Kannada for journalism, then the Kannada for a documentary, then say \emph{ōduga} again
\end{itemize}

\begin{tcolorbox}[breakable,colback=teal!4,colframe=teal!35!black,title={Before you move on}]
What does \kn{ಓದುಗ} mean? (\textquotedblleft{}A reader\textquotedblright{}.) Say it once more.
\end{tcolorbox}
\section[\texorpdfstring{\kn{ಕೇಳುಗ} (kēḷuga)}{ಕೇಳುಗ (kēḷuga)}]{\kn{ಕೇಳುಗ} (kēḷuga) — a listener}

\label{lesson:KA-C370-keluga}



\begin{quote}
Before the new one: say the Kannada for a documentary, then the Kannada for a reader.
\end{quote}

\subsection*{You'll want to know: \kn{ಕೇಳುಗ}}
\textbf{\kn{ಕೇಳುಗ}} — \emph{kēḷuga} — \textquotedblleft{}a listener\textquotedblright{}.

\begin{quote}
\kn{ರೇಡಿಯೋ} \kn{ಕಾರ್ಯಕ್ರಮಕ್ಕೆ} \kn{ತುಂಬಾ} \kn{ಕೇಳುಗರು} \kn{ಫೋನ್} \kn{ಮಾಡಿದರು}.
\end{quote}

\begin{quote}
\emph{rēḍiyō kāryakramakke tumbā kēḷugaru phōn māḍidaru.}
\end{quote}

\begin{quote}
Many listeners phoned the radio programme.
\end{quote}

Plural: \kn{ಕೇಳುಗರು}, \emph{kēḷugaru}.

Like \kn{ಓದುಗ}: \kn{ಕೇಳು}, to listen, with -\kn{ಗ}.

\begin{cousinweb}
\textbf{\kn{ಕೇಳುಗ}} (\emph{kēḷuga}) is a native Kannada noun: \emph{kēḷu}, \textquotedblleft{}to hear, to listen\textquotedblright{} (DEDR 2017), with the ending -\emph{ga}, \textquotedblleft{}one who does\textquotedblright{}.
\end{cousinweb}

\begin{grammarlens}[title={What you've built}]
One more word for telling what happened.
\end{grammarlens}

\subsection*{Your turn}
\begin{itemize}
\raggedright
  \item \emph{Say it:} \emph{kēḷuga}
  \item \emph{Say it:} \emph{kēḷuga}, once more
  \item \emph{Say it:} \emph{rēḍiyō kāryakramakke tumbā kēḷugaru phōn māḍidaru.} — Many listeners phoned the radio programme.
  \item \emph{Recall:} say the Kannada for a documentary, then the Kannada for a reader, then say \emph{kēḷuga} again
\end{itemize}

\begin{tcolorbox}[breakable,colback=teal!4,colframe=teal!35!black,title={Before you move on}]
What does \kn{ಕೇಳುಗ} mean? (\textquotedblleft{}A listener\textquotedblright{}.) Say it once more.
\end{tcolorbox}
\section[\texorpdfstring{\kn{ಕುಖ್ಯಾತ} (kukhyāta)}{ಕುಖ್ಯಾತ (kukhyāta)}]{\kn{ಕುಖ್ಯಾತ} (kukhyāta) — notorious}

\label{lesson:KA-C370-kukhyata}



\begin{quote}
Before the new one: say the Kannada for a reader, then the Kannada for a listener.
\end{quote}

\subsection*{You'll want to know: \kn{ಕುಖ್ಯಾತ}}
\textbf{\kn{ಕುಖ್ಯಾತ}} — \emph{kukhyāta} — \textquotedblleft{}notorious\textquotedblright{}.

\begin{quote}
\kn{ಆ} \kn{ಕುಖ್ಯಾತ} \kn{ಕಳ್ಳನನ್ನು} \kn{ಕೊನೆಗೆ} \kn{ಪೊಲೀಸರು} \kn{ಹಿಡಿದರು}.
\end{quote}

\begin{quote}
\emph{ā kukhyāta kaḷḷanannu konege polīsaru hiḍidaru.}
\end{quote}

\begin{quote}
The police finally caught that notorious thief.
\end{quote}

\kn{ಕು}-, bad, and \kn{ಖ್ಯಾತ}, known: well known for bad deeds. \kn{ಪ್ರಖ್ಯಾತ} is famous.

\begin{cousinweb}
\textbf{\kn{ಕುಖ್ಯಾತ}} (\emph{kukhyāta}) is borrowed from Sanskrit: \emph{ku-}, \textquotedblleft{}bad\textquotedblright{}, with \emph{khyāta}, \textquotedblleft{}named, well known, notorious\textquotedblright{} (Monier-Williams).
\end{cousinweb}

\begin{grammarlens}[title={What you've built}]
One more word for telling what happened.
\end{grammarlens}

\subsection*{Your turn}
\begin{itemize}
\raggedright
  \item \emph{Say it:} \emph{kukhyāta}
  \item \emph{Say it:} \emph{kukhyāta}, once more
  \item \emph{Say it:} \emph{ā kukhyāta kaḷḷanannu konege polīsaru hiḍidaru.} — The police finally caught that notorious thief.
  \item \emph{Recall:} say the Kannada for a reader, then the Kannada for a listener, then say \emph{kukhyāta} again
\end{itemize}

\begin{tcolorbox}[breakable,colback=teal!4,colframe=teal!35!black,title={Before you move on}]
What does \kn{ಕುಖ್ಯಾತ} mean? (\textquotedblleft{}Notorious\textquotedblright{}.) Say it once more.
\end{tcolorbox}
